\documentclass[11pt]{article}

\usepackage[utf8]{inputenc}
\usepackage[T1]{fontenc}
\usepackage{geometry}
\usepackage{amsmath}
\usepackage{amssymb}
\usepackage{booktabs}
\usepackage{hyperref}
\usepackage{xurl}
\usepackage{xcolor}
\usepackage{listings}
\usepackage{enumitem}
\usepackage{parskip}
\usepackage{graphicx}
\graphicspath{{../images/}}

\geometry{margin=1in}

\lstset{
  language=Java,
  basicstyle=\ttfamily\small,
  keywordstyle=\color{blue},
  commentstyle=\color{gray},
  stringstyle=\color{teal},
  showstringspaces=false,
  columns=fullflexible,
  frame=single,
  framerule=0.2pt,
  breaklines=true
}

\setlist[itemize]{topsep=0.3em,itemsep=0.2em}
\setlist[enumerate]{topsep=0.3em,itemsep=0.2em}

% Simple per-section source line.
\newcommand{\notesources}[1]{\par\noindent{\small\color{gray}\textit{Sources:} \sloppy #1\par}}

% Unit-wise source strings for this subject (used by slides/notes).
% Path is relative to the lecture `latex/` folder (the compilation CWD).
% OOPS with Java (CSEG1044) — unit-wise sources
%
% These are short, student-facing reference strings intended to be used with:
%   - \slidesources{...}  (in beamer slides)
%   - \notesources{...}   (in lecture notes)
%
% Style inspiration: other subjects in My-Subjects/ use semicolon-separated sources
% and include an "accessed" date for web documentation.

\newcommand{\OOPSUnitISources}{Schildt, \textit{Java: A Beginner's Guide} (10e), McGraw-Hill (Java fundamentals + OOP basics).; Oracle Java Tutorials: Getting Started + Language Basics + Classes and Objects (accessed \today).; Java SE API docs (e.g., \texttt{String}, \texttt{Scanner}) (accessed \today).}

\newcommand{\OOPSUnitIISources}{Schildt, \textit{Java: The Complete Reference} (12e), McGraw-Hill (classes, inheritance, packages, interfaces).; Bloch, \textit{Effective Java} (3e), Addison-Wesley, 2018 (design choices; interfaces vs abstract classes).; Oracle Java Tutorials: Classes and Objects + Interfaces + Packages (accessed \today).; Java SE API docs (e.g., \texttt{Comparable}, \texttt{Runnable}) (accessed \today).}

\newcommand{\OOPSUnitIIISources}{Schildt, \textit{Java: The Complete Reference} (12e) (nested/inner classes, exceptions, threads, I/O).; Oracle Java Tutorials: Nested Classes + Exceptions + Concurrency + Basic I/O (accessed \today).; Java SE API docs (e.g., \texttt{Thread}, \texttt{AutoCloseable}) (accessed \today).}

\newcommand{\OOPSUnitIVSources}{Schildt, \textit{Java: The Complete Reference} (12e) (generics, lambdas, Swing, JDBC).; Bloch, \textit{Effective Java} (3e) (generics + lambdas).; Oracle Java Tutorials: Generics + Lambda Expressions + Swing + JDBC Basics (accessed \today).; Java SE API docs (e.g., \texttt{java.util.function}, \texttt{javax.swing}, \texttt{java.sql}) (accessed \today).}

\newcommand{\OOPSUnitVSources}{Schildt, \textit{Java: The Complete Reference} (12e) (Collections Framework + wrapper classes).; Oracle Java docs: Collections Framework overview (accessed \today).; Java SE API docs (\texttt{java.util} collections; wrapper classes in \texttt{java.lang}) (accessed \today).}

\newcommand{\OOPSUnitVISources}{Git SCM: \textit{Pro Git} (2e) (version control workflow), accessed \today.; GitHub Docs (repositories, issues, pull requests), accessed \today.; Oracle Java Tutorials: Swing + JDBC (accessed \today).; JUnit 5 User Guide (accessed \today).; Apache Maven Project: Getting Started (accessed \today).; Gradle User Manual: Getting Started (accessed \today).; UML class diagrams: UML 2.x overview/spec (accessed \today).}

\newcommand{\OOPSMiscWhyInterfacesSources}{Schildt, \textit{Java: The Complete Reference} (12e) (interfaces).; Bloch, \textit{Effective Java} (3e) (prefer interfaces where appropriate).; Oracle Java Tutorials: Interfaces (accessed \today).; Java SE API docs (e.g., \texttt{Comparable}, \texttt{Runnable}) (accessed \today).}



\title{Object Oriented Programming (CSEG1044)\\Unit 03 -- Lecture 03 Notes\\Multithreading + Synchronization}
\begin{document}
\maketitle
\tableofcontents

\section{Lecture Snapshot}
\begin{itemize}
  \item \textbf{Unit focus:} Nested Classes, Exceptions, Multithreading, and I/O
  \item \textbf{Main new ideas:} Concurrency vs parallelism (quick intuition), Threads in Java: \texttt{Thread} and \texttt{Runnable}, Thread lifecycle (simplified)
  \item \textbf{Course thread:} Use Multithreading + Synchronization to make Campus Resource Manager safer, more modular, and easier to recover when something goes wrong.
\end{itemize}
\notesources{\OOPSUnitIIISources}

\section{Prerequisites}
You should already be comfortable with:
\begin{itemize}
  \item classes/objects and methods,
  \item loops and basic debugging,
  \item the idea of shared variables and side-effects.
\end{itemize}

\section{Learning Outcomes}
After this lecture, you should be able to:
\begin{itemize}
  \item Create threads using \texttt{Thread} and \texttt{Runnable}.
  \item Describe thread lifecycle and understand priority limitations.
  \item Identify race conditions and fix them with \texttt{synchronized}.
\end{itemize}

\section{Suggested Reading Order}
\begin{enumerate}
  \item Read the \textbf{Prerequisite Bridge} first and confirm each recalled idea.
  \item Study the central explanation in this order: \textit{Concurrency vs parallelism (quick intuition)} $\rightarrow$ \textit{Threads in Java: \texttt{Thread} and \texttt{Runnable}} $\rightarrow$ \textit{Thread lifecycle (simplified)} $\rightarrow$ \textit{Thread priorities (what they mean and what they do not)} $\rightarrow$ \textit{The real problem: shared mutable state}.
  \item Trace the worked example line by line before checking short answers.
  \item Attempt the practice section without looking at the solution immediately.
  \item Finish with the exit questions and further-study suggestions to confirm retention.
\end{enumerate}
\notesources{\OOPSUnitIIISources}

\section{Lecture Overview}
Multithreading allows a Java program to do multiple tasks at the same time (or appear to do so).
You need threads when you care about:
\begin{itemize}
  \item \textbf{responsiveness}: UI apps should not freeze while doing I/O or long work,
  \item \textbf{throughput}: servers handle many requests concurrently,
  \item \textbf{parallelism}: CPUs have multiple cores that can run tasks in parallel.
\end{itemize}
\textbf{Big warning:} threads make programs harder to reason about. Most bugs come from shared mutable data.
This lecture focuses on the core ideas required by the syllabus: creating threads, priorities, and synchronization.
\notesources{\OOPSUnitIIISources}

\section{Prerequisite Bridge}
Before reading the new material in depth, do a fast recall of the older ideas listed below.
\begin{itemize}
  \item \textbf{Runnable as a contract}: Threading builds on the same interface-based behavior style already seen earlier. Main reference: \texttt{Unit 02 / Lecture 05 slides/notes}.
  \item \textbf{Failure handling}: Concurrent tasks still need clear handling when something goes wrong. Main reference: \texttt{Unit 03 / Lecture 02 slides/notes}.
\end{itemize}
This lecture only revisits those ideas briefly so that the main explanation can stay focused on the new concept sequence.
\notesources{\OOPSUnitIIISources}

\section{Key Terms (Quick)}
\begin{itemize}
  \item \textbf{Concurrency}: managing multiple tasks that overlap in time.
  \item \textbf{Parallelism}: tasks literally running at the same time on different cores.
  \item \textbf{Thread}: lightweight unit of execution within a process.
  \item \textbf{Race condition}: outcome depends on timing/interleaving of threads.
  \item \textbf{Critical section}: code that must not be executed by multiple threads at once.
  \item \textbf{Lock/monitor}: mechanism used by \texttt{synchronized}.
\end{itemize}

\section{Detailed Notes}
\subsection{Concurrency vs parallelism (quick intuition)}
\begin{itemize}
  \item \textbf{Concurrency}: tasks overlap in time (switching between tasks).
  \item \textbf{Parallelism}: tasks run at the same time on different cores.
\end{itemize}
You can have concurrency on a single-core CPU (by switching), and you can have parallelism on multi-core CPUs.

\subsection{Threads in Java: \texttt{Thread} and \texttt{Runnable}}
Two common ways to create threads:
\begin{itemize}
  \item extend \texttt{Thread} and override \texttt{run()},
  \item implement \texttt{Runnable} and pass it to a \texttt{Thread} (preferred in many designs).
\end{itemize}

\paragraph{Most important rule: \texttt{start()} vs \texttt{run()}.}
\begin{itemize}
  \item \texttt{start()} creates a \textbf{new OS-level thread} and then calls \texttt{run()} on that new thread.
  \item \texttt{run()} called directly is just a normal method call (no new thread).
\end{itemize}

\paragraph{Thread naming helps debugging.}
\begin{lstlisting}
Thread t = new Thread(task);
t.setName("worker-1");
t.start();
\end{lstlisting}

\subsection{Thread lifecycle (simplified)}
\begin{itemize}
  \item \textbf{New}: thread object created, not started.
  \item \textbf{Runnable/Running}: eligible to run or currently running.
  \item \textbf{Blocked/Waiting}: sleeping, waiting for lock, or waiting for I/O.
  \item \textbf{Terminated}: finished \texttt{run()}.
\end{itemize}

\subsection{Thread priorities (what they mean and what they do not)}
\textbf{Priority is a hint} to the scheduler. It is not a guarantee.
On many systems it may have little visible effect.
Do not rely on priority for correctness.

\subsection{The real problem: shared mutable state}
If each thread works on its own local variables, life is easy.
Bugs start when threads share and modify the same object/variable.

\paragraph{Race condition.}
A race condition happens when the outcome depends on timing/interleaving.
The classic example is \texttt{value++}:
\begin{itemize}
  \item read \texttt{value},
  \item add 1,
  \item write back.
\end{itemize}
If two threads interleave those steps, an increment can be lost.

\subsection{Synchronization with \texttt{synchronized} (core tool)}
\texttt{synchronized} protects a \textbf{critical section} using a lock (monitor).
Only one thread at a time can execute synchronized code guarded by the same lock.

\paragraph{Two forms you must know.}
\begin{itemize}
  \item synchronized instance method: locks on \texttt{this}.
  \item synchronized block: locks on a specific object (\texttt{synchronized(lock)\{...\}}).
\end{itemize}

\paragraph{Lock scope tip.}
Keep the synchronized part as small as possible, but large enough to keep shared state consistent.

\subsection{Deadlock (common pitfall)}
Deadlock occurs when:
\begin{itemize}
  \item thread A holds lock 1 and waits for lock 2,
  \item thread B holds lock 2 and waits for lock 1.
\end{itemize}
\textbf{Prevention rule:} always acquire multiple locks in a consistent order.

\notesources{\OOPSUnitIIISources}

\section{Running Project Connection}
Use Multithreading + Synchronization to make Campus Resource Manager safer, more modular, and easier to recover when something goes wrong.
\begin{itemize}
  \item Use Campus Resource Manager as the running case study, or map the same idea to your own project domain if you are using another example.
  \item Ask where today's idea should live: model, service, DAO, UI, or team workflow.
  \item Use the lecture example as a smaller version of the same design choice you will make in the capstone.
\end{itemize}
\notesources{\OOPSUnitIIISources}

\section{Common Mistakes and Confusions}
\begin{itemize}
  \item Assuming threads execute in a fixed order every run.
  \item Treating thread priority as a correctness guarantee.
  \item Synchronizing everything blindly.
\end{itemize}
\notesources{\OOPSUnitIIISources}

\section{Worked Example (tutorial): see a race condition and fix it}
We will build a tiny counter and run two threads on it.

\subsection*{Step 1: Unsafe counter (race condition)}
\begin{lstlisting}
class Counter {
    private int value = 0;
    void inc() { value++; } // not thread-safe
    int get() { return value; }
}
\end{lstlisting}

\subsection*{Step 2: Safe counter using \texttt{synchronized}}
\begin{lstlisting}
class SafeCounter {
    private int value = 0;
    synchronized void inc() { value++; }
    synchronized int get() { return value; }
}
\end{lstlisting}

\subsection*{Step 3: Run two threads and wait using \texttt{join()}}
\begin{lstlisting}
public class Demo {
    public static void main(String[] args) throws InterruptedException {
        SafeCounter c = new SafeCounter();

        Runnable task = new Runnable() {
            @Override public void run() {
                for (int i = 0; i < 100000; i++) c.inc();
            }
        };

        Thread t1 = new Thread(task);
        Thread t2 = new Thread(task);
        t1.start(); t2.start();
        t1.join(); t2.join();

        System.out.println(c.get()); // expected 200000
    }
}
\end{lstlisting}
\textbf{Why \texttt{join()}?} It makes the main thread wait until worker threads finish.
Without join, main may print the result too early.

\subsection{Another worked example: bank account withdrawal (classic)}
This example shows why checks and updates must be atomic (together).

\subsection*{Unsafe version (can go negative or allow double-withdraw)}
\begin{lstlisting}
class BankAccount {
    private int balance = 100;

    void withdraw(int amount) {
        if (amount <= 0) throw new IllegalArgumentException("amount");

        if (balance >= amount) {
            // force a context switch to make the bug easier to see
            try { Thread.sleep(1); } catch (InterruptedException ignored) {}
            balance -= amount;
        }
    }

    int getBalance() { return balance; }
}
\end{lstlisting}
Two threads can both pass the check before either subtracts.

\subsection*{Safe version (synchronized critical section)}
\begin{lstlisting}
class SafeBankAccount {
    private int balance = 100;

    synchronized void withdraw(int amount) {
        if (amount <= 0) throw new IllegalArgumentException("amount");
        if (balance >= amount) balance -= amount;
    }

    synchronized int getBalance() { return balance; }
}
\end{lstlisting}

\section{Demo / Observation Task}
Use this block as a short reproducible demo or observation task.
\begin{itemize}
  \item Reproduce the lecture's main example around \textit{Concurrency vs parallelism (quick intuition)} once without looking at the solution first.
  \item Change one input, rule, or design choice in Campus Resource Manager and predict the result before checking it.
  \item Record one success case, one edge case, and one failure case so the idea becomes observable instead of purely theoretical.
\end{itemize}
\notesources{\OOPSUnitIIISources}

\section{Practice Check (with brief solutions)}
\begin{enumerate}
  \item What is the difference between concurrency and parallelism?\par
  \textbf{Answer:} concurrency is overlapping tasks; parallelism is simultaneous execution on multiple cores.

  \item Which method starts a new thread: \texttt{start()} or \texttt{run()}?\par
  \textbf{Answer:} \texttt{start()}.

  \item Why is \texttt{value++} unsafe without synchronization?\par
  \textbf{Answer:} it's a read-modify-write sequence; interleaving can lose updates.
\end{enumerate}

\section{Self-Study Practice (with hints)}
\begin{enumerate}
  \item Write a program that runs two threads printing numbers 1--10 with thread names.
  \textbf{Hint:} use \texttt{Thread.currentThread().getName()}.

  \item Modify the counter example: remove \texttt{synchronized} and run 5 times; record outputs.
  \textbf{Hint:} race bugs may not appear every time.

  \item Create a shared bank account and two threads withdrawing money; prevent negative balance using synchronization.
  \textbf{Hint:} the check and update must be inside the same synchronized block/method.

  \item Try changing thread priorities and observe (note that results may vary).
\end{enumerate}

\section{Summary (what to remember)}
\begin{itemize}
  \item Use \texttt{start()} to run code in a new thread; \texttt{run()} is just a method call.
  \item Use \texttt{join()} when one thread must wait for another to finish.
  \item Race conditions happen with shared mutable state.
  \item \texttt{synchronized} protects critical sections using a lock (monitor).
  \item Priority is not a correctness tool; it is only a hint.
\end{itemize}

\section{Further Study}
\begin{itemize}
  \item Revisit the prerequisite references if any recall bullet still feels weak.
  \item Rework the lecture example with one changed assumption, input, or design choice.
  \item Read the textbook and documentation references listed in the source line for a second explanation of the same topic.
  \item Apply the same idea once inside Campus Resource Manager so that the concept moves from theory to design practice.
\end{itemize}
\notesources{\OOPSUnitIIISources}

\section{Exit Questions (with sample answers)}
\textbf{Q1.} What is a race condition?\par
\textbf{Sample answer:} A race condition happens when multiple threads access shared mutable state without proper synchronization, so the outcome depends on timing/interleaving.

\vspace{0.6em}
\textbf{Q2.} Why does \texttt{synchronized} help?\par
\textbf{Sample answer:} It ensures only one thread at a time can execute the protected critical section, preventing lost updates and inconsistent states.

\vspace{0.6em}
\textbf{Q3.} What is the main reason for synchronization?\par
\textbf{Sample answer:} It protects shared mutable state from inconsistent updates when multiple threads use it.

\end{document}
